\section{Functions}
\label{sec:functions}

\subsection{Functions}

A function assigns exactly one output to every input.

Examples include

\[
x(t),
\]

\[
y(t),
\]

and

\[
z(t).
\]

Together they define the complete position vector

\[
\boxed{
	\mathbf r(t)
	=
	x(t)\hat{\mathbf x}
	+
	y(t)\hat{\mathbf y}
	+
	z(t)\hat{\mathbf z}.
}
\]

Every moving particle therefore traces a curve through space.

%%%%%%%%%%%%%%%%%%%%%%%%%%%%%%%%%%%%%%%%%%%%%%%%%%%%%%%%%%%%%%%%%%%%%%%%%%%%%%%%

\subsection{Why Calculus?}

Imagine observing an electron moving along a straight line.

At different moments we measure

\[
\begin{array}{c|c}
	t\;(\mathrm{s}) & x\;(\mathrm{m})\\
	\hline
	0 & 0\\
	1 & 2\\
	2 & 5\\
	3 & 9
\end{array}
\]

The particle is clearly moving.

More importantly, its speed is changing.

The natural question is

\begin{center}
	\emph{How can we describe continuous change mathematically?}
\end{center}

The answer is calculus.

Calculus was independently developed by Isaac Newton and Gottfried Wilhelm
Leibniz during the seventeenth century to describe continuously changing
physical systems.

%%%%%%%%%%%%%%%%%%%%%%%%%%%%%%%%%%%%%%%%%%%%%%%%%%%%%%%%%%%%%%%%%%%%%%%%%%%%%%%%
